\frenchspacing
\chapter{Introduction}
\epigraph{“If people knew how hard I had to work to gain my mastery, it would not seem so wonderful at all.”}{Michelangelo Buonarroti}

% Learning Techniques
\section{Learning Techniques}
The following is the five-step Feynman Technique:
\begin{enumerate}
    \item Choose the topic.
    \item Explain it as if you are explaining it to someone with no prior knowledge.
    \item Find where your explanation breaks down or lacks clarity.
    \item Review the material.
    \item Rewrite using language that is easily understood with clear examples.
\end{enumerate}
The following is my own personal process to write my conceptual understanding of the material I learned:

\begin{enumerate}
    \item \textbf{Take notes.}

    \item \textbf{Question.}
    \begin{itemize}
        \item What are we trying to learn?
        \item Why are we trying to learn this material?
        \item How will this apply to my overall goal?
    \end{itemize}

    \item \textbf{Extract.}
    \begin{itemize}
        \item Gather information from sources (reading, notes, slideshows, videos, etc.).
        \item Pull out information that will help us understand the answers to our questions.
    \end{itemize}

    \item \textbf{Load.}
    \begin{itemize}
        \item Internalize the information and find a way to connect it all together.
    \end{itemize}

    \item \textbf{Transform.}
    \begin{itemize}
        \item Synthesize and re-express the information in your own words.
        \item Use the Feynman Technique mentioned above to synthesize the information.
    \end{itemize}

    \item \textbf{Teach.}
    \begin{itemize}
        \item Using the Feynman style, try explaining it to a friend and find gaps.
        \item Review and refine.
    \end{itemize}
\end{enumerate}

\section{Pledge of Commitment}

\textit{Nihil Sine Labore} - Nothing without hard work or labor.